\documentclass[11pt]{article}

\usepackage[a4paper,margin=1in]{geometry}
\usepackage{amsmath,amssymb,amsthm,mathtools}
\usepackage{bm}
\usepackage{physics}
\usepackage{booktabs}
\usepackage{hyperref}

\hypersetup{
colorlinks=true,
linkcolor=blue,
citecolor=blue,
urlcolor=blue
}

\title{
Relational Emergence Model (REM 5.12)\\
Variational Selection of Tensor-Product Structure\\
in Quantum Systems
}

\author{
Yoshifumi Maruko\\
Independent Researcher\\
Yamagata, Japan
}

\date{2026}

\newtheorem{definition}{Definition}
\newtheorem{remark}{Remark}

\begin{document}

\maketitle

\begin{abstract}

Quantum theory does not uniquely determine subsystem structure.

We propose a variational principle on factorization space:

\[
F^* = \argmin_F \mathcal L(F)
\]

\[
\mathcal L = -\Phi_S + \lambda \mathcal C_H
\]

balancing informational articulation and dynamical separability.

The framework provides a generative extension of relational quantum mechanics,
in which subsystem structure itself is treated as a variational object.

\end{abstract}

\section{Problem}

Hilbert space admits multiple tensor-product structures:

\[
\mathcal H \cong \mathcal H_A \otimes \mathcal H_B
\]

Quantum theory alone does not specify which decomposition
corresponds to physically meaningful subsystems.

We treat subsystem structure as a variational problem
defined on factorization space.

\section{Three-layer architecture}

REM may be interpreted through three structural layers:

Layer 0:
Hilbert space without preferred factorization

Layer 1:
conceptual articulation of candidate bipartitions

Layer 2:
dynamical stabilization of selected structures

Layer 1 is represented by a conceptual map:

\[
D:
\mathcal H
\rightarrow
\{\mathcal H_A \otimes \mathcal H_B\}
\]

which specifies candidate subsystem structure.

D is not a physical operator but a structural articulation rule.

Effective realizations of D are represented by optimization on factorization space.

\begin{remark}

REM introduces a generative extension of relational quantum mechanics
by treating subsystem articulation itself as a variational object.

Relational structure is not assumed; it is selected.

\end{remark}

\section{Factorization space}

Let:

\[
N=n_A n_B
\]

Tensor-product structures form coset manifold:

\[
\mathcal F
\cong
U(N)/(U(n_A)\otimes U(n_B))
\]

Each point:

\[
F \in \mathcal F
\]

represents a candidate subsystem decomposition.

\section{Variational principle}

Optimal factorization:

\[
F^*=\argmin_F \mathcal L(F)
\]

objective functional:

\[
\mathcal L
=
-\Phi_S
+
\lambda \mathcal C_H
\]

λ controls tradeoff between informational articulation
and dynamical separability.

\section{Informational articulation}

We define:

\[
\Phi_S=I(A:B)
\]

mutual information:

\[
I(A:B)
=
S(\rho_A)+S(\rho_B)-S(\rho)
\]

von Neumann entropy:

\[
S(\rho)
=
-\mathrm{Tr}(\rho\log\rho)
\]

Large Φ_S corresponds to strongly articulated relational structure.

\section{Interaction projection}

For operator X:

\[
X_A
=
\frac{1}{n_B}\mathrm{Tr}_B(X)
\]

\[
X_B
=
\frac{1}{n_A}\mathrm{Tr}_A(X)
\]

interaction component:

\[
X_{\partial F}
=
X
-
(X_A\otimes I_B)
-
(I_A\otimes X_B)
\]

This Hilbert–Schmidt decomposition isolates cross-subsystem coupling.

\section{Dynamical separability}

Hamiltonian decomposition:

\[
H=H_A+H_B+H_{\partial F}
\]

interaction cost:

\[
\mathcal C_H
=
\mathrm{Tr}(\rho H_{\partial F}^2)
\]

Quadratic form ensures positive measure of coupling strength.

\section{Pull-back parametrization}

Fix reference factorization.

Represent candidate structures via unitary transformation:

\[
U(\theta)\in U(N)
\]

state:

\[
\rho(\theta)
=
U^\dagger\rho_0U
\]

Hamiltonian:

\[
H(\theta)
=
U^\dagger H_0 U
\]

Gauge redundancy:

\[
U(\theta)
\sim
U(\theta)(U_A\otimes U_B)
\]

Optimization therefore occurs on quotient manifold.

\section{Gradient structure}

density variation:

\[
\delta\rho
=
i[\rho,G_k]\delta\theta_k
\]

entropy variation:

\[
\delta S(\rho)
=
-\mathrm{Tr}(\delta\rho(\log\rho+I))
\]

For reduced states:

\[
\delta\Phi_S
=
\delta S(\rho_A)
+
\delta S(\rho_B)
-
\delta S(\rho)
\]

interaction variation:

\[
\delta \mathcal C_H
=
\mathrm{Tr}(\delta\rho H_{\partial F}^2)
+
\mathrm{Tr}(\rho\delta(H_{\partial F}^2))
\]

stationary condition:

\[
\nabla_F\mathcal L=0
\]

defines optimal tensor-product structure.

\section{Role of λ}

λ acts as external control parameter
determining relative weight of information and dynamics.

Possible interpretation:

\[
\lambda
\sim
\frac{J^2}{\Gamma_I}
\]

where J is interaction scale
and Γ_I is information production rate.

Precise microscopic derivation remains future work.

\section{Three-qubit illustration}

state:

\[
|\Psi\rangle
=
\sqrt{0.6}|000\rangle
+
\sqrt{0.3}|111\rangle
+
\sqrt{0.1}|101\rangle
\]

reduced density matrix for qubit A:

\[
\rho_A
=
\begin{pmatrix}
0.6+0.3 & 0 \\
0 & 0.1
\end{pmatrix}
=
\begin{pmatrix}
0.9 & 0 \\
0 & 0.1
\end{pmatrix}
\]

entropy:

\[
S(A)=0.469
\]

mutual information:

\[
I(A:BC)=1.942
\]

interaction costs differ across factorizations,
producing regime competition:

information-dominated regime selects correlated structure

dynamics-dominated regime selects separable structure

crossover:

\[
\lambda^*\approx0.28
\]

\section{Relation to prior work}

Carroll and Singh identify preferred tensor-product structure
through dynamical robustness criteria,
minimizing interaction leakage and entanglement growth.

REM generalizes this approach by introducing explicit informational objective Φ_S,
leading to hybrid informational–dynamical optimization on factorization space.

Zanardi defines subsystem structure operationally through observable algebra.

REM instead treats tensor-product structure as variational degree of freedom.

Loizeau and Sels infer subsystem structure from spectral properties alone.

REM uses both state information and Hamiltonian structure.

\begin{center}
\begin{tabular}{lccc}
\toprule
framework & information & dynamics & generative \\
\midrule
Zanardi & implicit & yes & no \\
Carroll-Singh & no & yes & no \\
Loizeau-Sels & no & indirect & no \\
Zurek & implicit & yes & dynamic \\
REM & yes & yes & yes \\
\bottomrule
\end{tabular}
\end{center}

Generative refers to treating subsystem structure itself
as an optimization variable.

\section{Open-system perspective}

Layer 2 corresponds to stability of selected structure
under system dynamics.

In open quantum systems,
environmental monitoring may favor stable pointer structures,
as described in Zurek's einselection framework.

This suggests possible extension:

\[
\mathcal C_E
=
\sum_k
\mathrm{Tr}
(
\rho
L_k^\dagger L_k
)_{\partial F}
\]

linking REM to decoherence-induced stability.

A full dissipative treatment is left for future work.

\section{Multipartite extension}

General structure:

\[
\mathcal H
=
\bigotimes_i \mathcal H_i
\]

total correlation:

\[
T(\rho)
=
\sum_i S(\rho_i)-S(\rho)
\]

supports hierarchical emergence of subsystem structure.

\section{Computational considerations}

dimension:

\[
\dim(\mathcal F)
=
N^2-n_A^2-n_B^2
\]

possible strategies:

local Hamiltonian constraints

tensor-network parametrization

hierarchical coarse graining

\section{Conclusion}

Subsystem structure may be selected variationally:

\[
F^*
=
\argmin_F
\mathcal L(F)
\]

REM provides generative extension of relational quantum mechanics,
formulating subsystem emergence as optimization on factorization space.

\begin{thebibliography}{9}
\bibitem{rovelli} C Rovelli. Relational quantum mechanics. Int J Theor Phys 35 1637 (1996)
\bibitem{zanardi} P Zanardi. Virtual subsystems. Phys Rev Lett 92 060402 (2004)
\bibitem{carroll} S Carroll A Singh. Quantum Mereology. Phys Rev A 103 022213 (2021)
\bibitem{loizeau} N Loizeau D Sels. Subsystem extraction from spectrum. Found Phys 55 (2025)
\bibitem{zurek} W Zurek. Decoherence and einselection. Rev Mod Phys 75 715 (2003)
\bibitem{frauchiger} D Frauchiger R Renner. Nat Commun 9 3711 (2018)
\end{thebibliography}

\end{document}
